\scanpage{38}
% L1
Si $\Sigma^+(X)=$ ens des classes d'objets simples de $M^+(X)$, $\Sigma(X)=\underset{\substack{\text{par}\\e\mapsto\xi e}}{\underrightarrow{\lim}}\Sigma^+(X)$, on aura \deleted{\unsure{pour}} \deleted{$M(K)$} un diagramme comm.
\diagram{\begin{tikzcd}[ampersand replacement=\&,column sep=small,row sep=large]
M^+(X)\arrow[d,hook] \& = \& \mathbb Z^{(\Sigma^+(X))}\arrow[d,hook]\\
M(X) \& = \& \mathbb Z^{(\Sigma(X)}
\end{tikzcd}}

% L6
Alors $M_i^{\unsure{+}}(X)$ est défini en termes d'une partie $\Sigma^+_{(i)}(X)$ de $\Sigma_i(X)$, d'où
\[
M^+_i(X)=\mathbb Z^{\Sigma^+_{(i)}(X)}
\]
ce qui donne une filtration croissante sur l'anneau $M^+(X)$, et une filtration analogue sur $M(X)$.

% L12
En fait, posant
\[
\Sigma_i^+(X)=\Sigma^+_{(i)}(X)-\Sigma^+_{(\unsure{i-1|i+1})}(X)
\]
on aura $=$
\[
\Sigma_i^+(X)\cdot\Sigma_j^+(X)\subset\mathbb Z^{\Sigma^+_{\unsure{i+j|i-j}}(X)}
\]
donc $M^+(X)$ devient un \uline{anneau}
% archive sheet number: 37
